\documentclass[12pt]{article}
\usepackage[utf8]{inputenc}
\usepackage[T1]{fontenc}
\usepackage[margin=1in]{geometry}
\usepackage{setspace}
\usepackage{xcolor}
\usepackage[colorlinks=true, urlcolor=blue!60!black, linkcolor=blue!60!black]{hyperref}
\setlength{\parindent}{2em}
\setlength{\parskip}{1.25em}
\renewcommand{\baselinestretch}{1.0}
\begin{document}
\begin{center}
{\LARGE \textit{Timbs v.\ Indiana}, 586 U.S.\ 146 (2019)\par}
\vspace{0.8em}
{\large Case Brief --- Four-Page Edition\par}
\end{center}
\vspace{1em}

\section{Memory Jogger}

The Eighth Amendment's Excessive Fines Clause is incorporated against the States, so Indiana could not escape excessive-fines scrutiny when seeking civil \textit{in rem} forfeiture of a \$42,000 Land Rover for a drug offense carrying a \$10,000 maximum fine (\href{https://www.courtlistener.com/opinion/4591916/timbs-v-indiana/}{CourtListener}).

\section{Detailed Case Facts}

Tyson Timbs used about \$42,000 of life insurance proceeds from his father's death to buy a 2012 Land Rover LR2. In May 2013, struggling with opioid addiction, he drove it to sell about four grams of heroin to undercover officers in Grant County, Indiana. Police arrested him and seized the vehicle. Timbs pleaded guilty to dealing in a controlled substance and conspiracy to commit theft, and received one year of home detention, five years of probation with addiction treatment, and \$1,203 in fees and costs. Indiana hired a private law firm to bring a civil \textit{in rem} forfeiture action against the Land Rover. The vehicle was worth more than four times the \$10,000 maximum fine for his offense.

\section{Procedural History}

The trial court denied forfeiture as ``grossly disproportionate to the gravity of Timbs's offense'' under the Excessive Fines Clause, and the Court of Appeals of Indiana affirmed. The Indiana Supreme Court reversed without deciding excessiveness, holding that the Clause had never been incorporated and bound only the federal government. Represented by the Institute for Justice, Timbs obtained certiorari; the Court vacated and remanded on February 20, 2019. On remand, the Indiana Supreme Court held the forfeiture unconstitutionally excessive.

\section{Judicial Votes}

Vacated and remanded, \textbf{9--0}.
\begin{itemize}
\item \textbf{Majority:} Justice Ginsburg, joined by Chief Justice Roberts and Justices Breyer, Alito, Sotomayor, Kagan, Gorsuch, and Kavanaugh.
\item \textbf{Concurrence:} Justice Gorsuch.
\item \textbf{Concurring in the judgment:} Justice Thomas.
\end{itemize}

\section{Holding}

The Excessive Fines Clause is incorporated through the Due Process Clause because protection against excessive punitive economic sanctions is ``fundamental to our scheme of ordered liberty'' and ``deeply rooted in this Nation's history and tradition.'' Once incorporated, a guarantee applies identically to the States, including to partly punitive civil \textit{in rem} forfeitures under \textit{Austin v.\ United States}. The incorporation question asks whether the right itself, not each application, is fundamental. The Court did not decide whether this forfeiture was excessive, set an excessiveness test, or decide whether ability to pay matters.

\section{Analysis of Opinions}

\textbf{Majority (Ginsburg).} Applying \textit{McDonald v.\ City of Chicago}, Ginsburg traced the Clause's ``venerable lineage'' from Magna Carta, which required amercements proportioned to the offense and sparing a livelihood, through the English Bill of Rights of 1689, the Virginia Declaration of Rights, and state constitutions in 1787 (8 of 13), 1868 (35 of 37), and today (all 50). She stressed the Black Codes' use of vagrancy fines to coerce freedmen into labor, and that fines, unlike imprisonment, raise revenue and can chill dissent. Indiana's request to overrule \textit{Austin} was not properly presented. Its argument that \textit{in rem} forfeiture itself must be deeply rooted misread the doctrine: the First Amendment reaches social media (\textit{Packingham}) and the Fourth reaches cell phones (\textit{Riley}) without separate incorporation.

\textbf{Gorsuch (concurring).} As an original matter, the Privileges or Immunities Clause may be the right vehicle, but nothing turned on it.

\textbf{Thomas (concurring in the judgment).} Thomas rejected ``the oxymoronic `substantive' `due process' doctrine'' as a legal fiction and, after a historical survey from Henry I's Charter of Liberties to the 39th Congress's debates on Black Code fines, held freedom from excessive fines a privilege of American citizenship.

\section{Examples: Future Applications}

\begin{enumerate}
\item \textbf{Same-Side: Forfeiting a Home or Plane.} Seeking civil forfeiture of a \$150,000 plane or home for a misdemeanor with a \$2,500 maximum fine must be tested for gross disproportionality; the ``civil'' label does not matter.
\item \textbf{Same-Side: Code-Enforcement Fines.} Daily fines totaling \$160,000 for overgrown grass are constrained by the incorporated Clause, which bars punitive tolls as a substitute for taxes.
\item \textbf{Opposite-Side: Remedial Recoveries.} Compensatory damages or back taxes with standard interest are not extracted ``as punishment for some offense.''
\item \textbf{Opposite-Side: Contraband and Proceeds.} Forfeiting \$50,000 in proven drug proceeds or illegal automatic weapons is generally treated as remedial.
\item \textbf{Fence-Sitter: Ability to Pay.} A \$1,500 fine that costs an indigent driver his license and livelihood invokes Magna Carta's livelihood principle, but the Court left ability to pay open, and lower courts divide.
\end{enumerate}

\section{Critique}

\textbf{Originalist.} Originalists and libertarians welcomed the check on civil forfeiture but, with Thomas and Gorsuch, faulted reliance on substantive due process rather than the Privileges or Immunities Clause; some note that \textit{Austin}'s validity remains open to originalist challenge.

\textbf{Progressive and doctrinal.} Scholars praised the attention to Black Code fines and modern revenue-driven fines and fees, but Beth Colgan and the \textit{Harvard Law Review} warn that \textit{Timbs} may prove symbolic unless courts consider ability to pay; the lack of a uniform excessiveness standard has fragmented state tests.

\section{Key Quotations}

\begin{quote}
``Incorporated Bill of Rights guarantees are `enforced against the States under the Fourteenth Amendment according to the same standards that protect those personal rights against federal encroachment.' Thus, if a Bill of Rights protection is incorporated, there is no daylight between the federal and state conduct it prohibits or requires.''
\end{quote}
\noindent --- Justice Ginsburg, for the Court, 586 U.S.\ at 150 (quoting \textit{McDonald v.\ Chicago}, 561 U.S.\ at 765).

\begin{quote}
``For good reason, the protection against excessive fines has been a constant shield throughout Anglo-American history: Exorbitant tolls undermine other constitutional liberties.''
\end{quote}
\noindent --- Justice Ginsburg, for the Court, 586 U.S.\ at 153--54.

\begin{quote}
``In considering whether the Fourteenth Amendment incorporates a protection contained in the Bill of Rights, we ask whether the right guaranteed---not each and every particular application of that right---is fundamental or deeply rooted.''
\end{quote}
\noindent --- Justice Ginsburg, for the Court, 586 U.S.\ at 156.

\end{document}
